\section{Theodicy and the Other Woman}

\begin{quotex}
Tis better to have loved and lost than never to have loved at all.
\flright{\textsc{Alfred Lord Tennyson}}

God made us to show forth His goodness and to share with us His everlasting happiness in heaven. \flright{\emph{Baltimore Catechism \#3}}
\end{quotex}


According to Scripture, the world created by God is ``good''. So why is there a question of evil? That is because the current world is not created by God, but by man. It is the world constituted by the knowledge of good and evil and hence death.

By the Principle of Plenitude\footnote{\url{https://www.gornahoor.net/?p=7668}}, confirmed in the Catechism, existence itself is good. This may be counter-intuitive because we are accustomed to associate the good with sensual pleasure and nice emotions. Hence, pain, unsatisfied desires, and suffering must be evil. By this standard, then, life itself is suffering, since there is always desire. Keep in mind that Buddha, as Josaphat\footnote{\url{https://www.gornahoor.net/?p=7591}}, is considered a saint in the Christian church. Nevertheless, in the Medieval Tradition, the good of being surpasses the suffering of living.

The natural world is held in place by hunger, sex, and fear. The human race is motivated by the desire for things, the entanglements arising from following the sex instincts, and the reaction to fears of all types. This is borne out mightily by Hollywood films. It is an easy exercise to classify films based on their plots.

\begin{itemize}


\item Does the plot emphasize greed and conspicuous consumption?

\item Is it about love and sex and all that derives from them?

\item Does it reinforce common fears, or sometimes supernatural fears?
\end{itemize}


Of course, these motivating factors get resolved in the end, not by transcending them from a higher plane, but always on the natural plane.

\paragraph{The Other Woman}


In the Genesis story, creation is good until it is noticed that man is alone. Therefore, Eve is good for him and necessarily comes into being. Unfortunately, in the human world, the relationship between a man and a woman is as much a cause of pain and suffering as it is of joy. This duality is explored in the recent film, The Other Woman\footnote{\url{http://www.imdb.com/title/tt2203939/}}.

Kate is married to the handsome, wealthy, and successful Mark, assuming all is well with their marriage. Kate is scatter-brained, yet pops out ideas not unlike the way a flatulent person produces farts. The husband allegedly gains financial success from them, although not legitimately, but only fraudulently.

Meanwhile, his mistress Carly, a Manhattan lawyer, tries to surprise him at his house in Connecticut, but is surprised herself to find his wife at home. Passing over the irrelevant details, the two women, each feeling betrayed, form an unlikely bond. Eventually, a second mistress shows up, ``acted'' by the model Kate Upton as Amber (of course). Although not strictly necessary for the plot, she serves another purpose to be explained below.

\paragraph{Detecting Lies}


So this brings up the question: although Kate and Carly had such high hopes for their relationship with Mark, why were they unable to recognize his obvious lies? The glib answer is that everyone lies, which is certainly true. Nevertheless, they won't be visible to someone who lies to herself.

Several years ago I was called to jury duty. In the voir dire, the defendant's attorney asked all the potential jurors how could they tell that someone was lying. Virtually every woman responded that she could tell by looking in his eyes. Now even a modicum of real life experience should have disabused them of that opinion. I know as a fact that many women are quite happy to accept a lie.

At my turn, I responded with something stupid, along these lines: I would consider the plausibility of the testimony, its consistency both with the witness and other witnesses, and so on. For that, I was immediately dismissed and the women stayed behind. So I never learned of the outcome.

\paragraph{Chemistry}


While a man should be looking for a helpmate suitable for him, women have other plans. They may claim they want someone intelligent, funny, successful, interesting, and so on, but ultimately, in most cases today, it comes down to what is called ``chemistry''. That is, rather than evaluating a potential partner intellectually, they are looking for a biochemical event to make the decision for them. Obviously, the three women had ``chemistry'' at some point with Mark, for all the good it did them. Nevertheless, at the end when they each move on, they still rely on chemistry.

Now a man, too, is often overcome with his own chemistry, but the goal in that case is seldom a permanent relationship. That is why the aristocrat of the soul must learn to see circumstances rightly with his third eye rather than with his one-eyed member.

\paragraph{Amber's Ass}


Kate Hudson's ass is on display several times ... in close up ... covered with a bikini. Apparently the ass is replacing the breast as the primary object of man's lust. Ass play, as normative, puts heterosexuality and homosexuality on an equal basis. For example, the new season of \emph{Girls} on HBO, hoping apparently to challenge any lingering Puritanism in the USA, featured a totally gratuitous scene involving anilingus. It had no relationship to the story and, since it was done over the kitchen sink, it was hardly a romantic moment.

Unfortunately, Lena Dunham must have missed Two Girls and a Guy\footnote{\url{https://en.wikipedia.org/wiki/Two_Girls_and_a_Guy}}, released 15 years ago, which featured a similar scene involving Heather Graham as the object of affection. At least in that case, the temptation was believable.

\paragraph{A Woman and her Dog}


Of course, the jilted wife had a dog, a large dog, as a symbol of unconditional love. When the dog crapped Carly's apartment, it was easily dismissed as a ``mistake''. However, when the husband confessed to a ``mistake'', he received no sympathy from any quarter. Perhaps if Kate had treated her husband like her dog, things may have been different. Specifically, she would have his meals ready on time, pick up all his crap, and scratch his belly at night. The fact that that probably sounds oppressive to most people when applied to a man, but not to a dog, reveals something ... I'm sure of it.

\paragraph{The Happy Ending}


Although this movie was panned by the critics, it was a financial success. So it is worthwhile to unravel the hidden assumptions in the script to understand its appeal.

The husband Mark is handsome, appealing to women, and a financial success. This he gained his advantages through deception and fraud, the only explanation for his position is ``white male privilege''. The three women are completely innocent victims, since all they wanted was to love and please Mark. Individually, they were powerless against him. Yet, in Nietzschean fashion, when they worked together as a group, they were able to topple Mark's position of dominance. In the end, Mark gets his comeuppance, losing his wife, his lovers, his money, and his career.

The film then appeals to two conflicting drives in contemporary women: the desire for independence versus the age-old propensity for female hypergamy. Carly gets Kate's brother along with his beach house in the Hamptons. Amber marries Carly's wealthy and much older father, so they can travel without financial concerns.

Kate, on the other hand, now free of Mark's patriarchal oppression, recapitulates a theme from The Purge\footnote{\url{https://www.gornahoor.net/?p=7412\#purge}}. She suddenly becomes a ``strong and independent'' woman, becoming successful in her own right from all her own ideas.

I can't recommend this movie, but it you must, then fast forward to the scenes with Olivia Culpo\footnote{\url{http://top-beautiful-women.com/usa/item/31-olivia-culpo}} if you just want to see an attractive woman.

\url{https://www.youtube.com/watch?v=uzTLRdWKah8} 

\flrightit{Posted on 2015-01-28 by Cologero}

\begin{center}* * *\end{center}

\begin{footnotesize}
\begin{sffamily}
\texttt{Matt on 2015-01-28 at 10:03 said:}

``At my turn, I responded with something stupid, along these lines: I would consider the plausibility of the testimony, its consistency both with the witness and other witnesses, and so on. For that, I was immediately dismissed and the women stayed behind.''

A good example showing how justice is not the organizing principle of the current legal system, and has instead been replaced by cunning and manipulation to get to the goal of ``winning''.

At a prior period in the histroy of the West, people had another term for ``chemistry'': lust.

\hfill

\texttt{Justin on 2015-01-28 at 13:01 said:}

``Apparently, in our homosexualized culture the ass is replacing the breast as the primary object of man's lust. Ass play, as normative, puts heterosexuality and homosexuality on an equal basis.''

I would say that there is a significant distinction to be made here between the buttocks and the rectum as objects of male attraction/lust.

\hfill

\texttt{Paulo on 2015-01-28 at 15:36 said:}

Regarding the hypergamy Cologero,it is also completely disturbed,the case that make modern women fell atracktion to a man,is demonic !!!

\hfill

\texttt{Paulo on 2015-01-28 at 15:39 said:}

Yes it is Justin,but people nowaday cannot make this distinction any more!

\hfill

\texttt{August on 2015-01-29 at 09:26 said:}

``At my turn, I responded with something stupid... ''

Such `stupid' responses are a great way to get out of all sorts of obligations in modern life, not just jury duty, so if anyone is looking for more free time, use them liberally in the right places.

And once a man is no longer regarding circumstances with his `one-eyed member', directly or indirectly, he might conclude that there has never been an easier place to live in than the modern world. Of course, one's temptation to stick his tongue in dark places may diminish, and that could cost him some social cachet, but nothing is free.

\hfill

\texttt{Logres on 2015-01-31 at 00:54 said:}

My father got out of jury duty by saying that he would use as his standard the book of Deuteronomy, which was God's Law. In his case, he was serious, although he probably knew and didn't mind the outcome.

\hfill

\texttt{Dan on 2015-01-31 at 14:47 said:}

I would really love to read your take on the movie ``Nightcrawler''. Personally I've never seen evil presented in such a powerful manner and the film's protagonist embodies the culture of materialism, narcissism and greed in it most horrifying manifestations. The film is actually full of symbolism such as the occurrence of the number 6 in may scenes, suggesting that this is actually what the embodiment of an Antichrist figure would be like.

\hfill

\texttt{ja on 2015-02-03 at 15:40 said:}

Oh how superior the Eye of Horus is to the Mouth of Isis !

\end{sffamily}
\end{footnotesize}

